% !TeX root = ../../Book3.tex
% 纲要标题：### G.1 半代数集合与量词消去

\section{半代数集合与量词消去}
\label{b3:sec:G01}

实零点会随系数的符号变化，仅有多项式等式不足以描述参数条件。把等式与不等式一起保存在半代数集合中，消去变量后仍可得到有限的符号描述。

% 纲要细目（与计划纲要同步）：
% * 实闭域上由多项式等式与不等式及有限布尔运算定义的集合；半代数映射
% * Tarski–Seidenberg 投影定理与实闭域量词消去，说明它们为何超出普通 Zariski 消去
% * 手算一元或二元投影例子；完整柱状代数分解算法不放入正文
%
% #### G.1.1 半代数集合与半代数映射
% #### G.1.2 Tarski–Seidenberg 投影定理与量词消去
% #### G.1.3 低维投影算例与柱状代数分解


% 正文按纲要完成；所列定理给出证明或证明概要。

\subsection{半代数集合与半代数映射}
\label{b3:subsec:G01:01}

实系数方程 \(x^2+1=0\) 没有实解，而 \(x^2-a=0\) 是否有实解取决于参数 \(a\) 的符号。这说明研究实零点时，等式与不等式必须一起考虑。本附录用 \(R\) 表示实闭域；涉及紧性、光滑性和测度时，另明确限定为 \(\mathbb R\)。第一卷附录 E 已介绍实闭域及多项式的介值性质。

\begin{definition}\label{b3:def:semialgebraic-set}
设 \(R\) 为实闭域，\(m,n\) 为非负整数。若 \(R^n\) 的一个子集可由有限多个条件 \(f(x)=0\)、\(g(x)>0\)（其中 \(f,g\in R[x_1,\ldots,x_n]\)）经有限次并、交、补得到，则称它为\term{半代数集合}[semialgebraic set]。若 \(A\subseteq R^m\)、\(B\subseteq R^n\) 都是半代数集合，且映射 \(\varphi:A\to B\) 的图为 \(R^{m+n}\) 的半代数子集，则称 \(\varphi\) 为\term{半代数映射}[semialgebraic map]。
\end{definition}

条件 \(g\ge0\) 等价于 \(g>0\) 或 \(g=0\)，故定义允许弱不等式。把有限个多项式的符号取值逐一列出，还可把任意半代数集合写成有限个由等式与严格不等式定义的集合之并。多项式映射是半代数映射；有理函数只在其分母非零的定义域上给出半代数映射。例如 \(x\mapsto1/x\) 的图像由 \(xy=1\) 定义。一般半代数映射不必连续：符号函数就是例子。上述用语可参阅 Mangolte《Real Algebraic Varieties》\cite[Appendix B, Definitions B.2.1--B.2.2]{Mangolte2020RealAlgebraicVarieties}。

\subsection{Tarski–Seidenberg 投影定理与量词消去}
\label{b3:subsec:G01:02}

\begin{theorem}[Tarski--Seidenberg 投影定理]\label{b3:thm:real-projection}
设 \(R\) 为实闭域，\(S\subseteq R^{n+m}\) 为半代数集合。坐标投影 \(\pi:R^{n+m}\to R^n\) 的像 \(\pi(S)\) 仍为半代数集合。
\end{theorem}

\begin{proof}
逐个投影，归约为消去一个变量 \(T\)。将定义 \(S\) 的布尔公式中出现的多项式记为
\[
p_1(x,T),\ldots,p_l(x,T).
\]
对一个固定参数 \(x\)，我们要判定哪几个符号元组能由某个 \(T\in R\) 实现；若这一判定只需有限个参数多项式的符号，投影便为半代数集。

先按各 \(p_i\) 的 \(T\) 次数分支，次数由系数的零与非零条件确定；恒零多项式单独处理。用 Euclid 算法从其余 \(p_i\) 的乘积取得首一平方自由部分 \(P\)。每次除法按当前首项系数是否为零分支，次数下降保证只有有限步、有限个分支；各分支上 \(P\) 的系数是参数的有理函数，所用分母已限定非零。若 \(P=1\)，所有 \(p_i\) 在整条直线上符号恒定，直接由首项系数判定。以下设 \(\deg P=d>0\)。

令 \(B=R[T]/(P)\)。对于任意 \(q\in R[T]\)，在基 \(1,T,\ldots,T^{d-1}\) 上计算迹二次型
\[
H_q(u,v)=\operatorname{Tr}_{B/R}(quv).
\]
它的矩阵只用 \(P,q\) 的系数作有限次域运算。实闭域上，平方自由代数 \(B\) 分解成若干 \(R\) 与 \(R(\mathrm i)\) 的直积。在实根 \(a\) 对应的分量中，二次型为 \(q(a)u^2\)，符号差为 \(\operatorname{sign}q(a)\)；在非实共轭对对应的分量中，非零迹形式有一正一负，\(q\) 为零时形式也为零。因此
\[
\operatorname{signature}(H_q)
=\sum_{P(a)=0,\ a\in R}\operatorname{sign}q(a).
\]
这里符号差指正惯性指数减负惯性指数。对称矩阵可按非零对角元作合同消元；若全部对角元为零但某个非对角元非零，就取对应的二阶块作为枢轴。按枢轴的零与符号分支，有限步可算出惯性指数。所以迹形式的符号差由有限多个参数有理函数的符号决定，而分母非零时这些仍能写成多项式符号条件。迹形式的这一根计数原理也见\cite[Section 2.4, Theorem 2.8]{Sturmfels2002PolynomialEquations}。

把所有 \(p_i\) 及其 \(T\) 导数列为 \(q_1,\ldots,q_b\)。对每个 \(e\in\{0,1,2\}^b\)，计算 \(H_{\prod_jq_j^{e_j}}\) 的符号差。若 \(N_\sigma\) 表示这些 \(q_j\) 在 \(P\) 的实根上取符号元组 \(\sigma\in\{-1,0,1\}^b\) 的根数，则这些差给出有限线性方程
\[
Q_e=\sum_\sigma N_\sigma\prod_j\sigma_j^{e_j},
\qquad 0^0=1.
\]
一变量的系数矩阵是 \(1,s,s^2\) 在 \(s=-1,0,1\) 上的取值矩阵，行列式非零；多变量矩阵是它们的张量积，故也可逆。于是知道各 \(Q_e\) 就能确定每个 \(N_\sigma\)，从而判定哪些根上的符号元组实际出现。

这些数据也确定各根之间开区间上的符号元组。在一个根 \(a\) 的右边，\(p_i\) 的符号由其在 \(a\) 处首个非零导数的符号给出；这个符号在到下一根之前不变。最左边区间的符号由 \((-1)^{\deg p_i}\) 乘首项系数的符号给出。有限实根之间均有元素，实闭域的多项式介值性质保证开区间内没有未计入的变号。因此根上的元组、各根右侧的元组与最左侧元组已经穷尽所有 \(T\) 的可能性，无须先求出根的位置。

检查这些有限元组是否满足原布尔公式，就判定了 \(\exists T\)。前述所有分支均由有限个参数多项式的零、正、负条件给出；合并满足条件的分支，得到投影的半代数描述。反复消去 \(m\) 个变量即得结论。
\end{proof}

投影性质的概述见\cite[Chapter 2, p.~23]{BochnakCosteRoy1998RealAG}。条件
\[
\exists y_1,\ldots,y_m\quad \Phi(x,y)
\]
描述的正是投影，其中 \(\Phi\) 是多项式符号条件的有限布尔组合。定理把它改写成只含 \(x\) 的有限布尔组合；全称量词则用 \(\forall y\,\Phi\iff\neg\exists y\,\neg\Phi\) 处理。逐层消去量词，便得到实闭域理论的量词消去。

\ProseParagraphBreak
这与\cref{b3:thm:elimination-projection-closure}的任务不同。该定理在代数闭域上用消去理想描述投影的 Zariski 闭包，不能单凭等式记住实数的符号。例如 \(I=(y^2-x)\subseteq R[x,y]\) 满足 \(I\cap R[x]=0\)，但实投影是 \(\{x\ge0\}\)，并非整个 \(R\)。实代数消去要保留的是实际投影，因此输出自然包含不等式。

\subsection{低维投影算例与柱状代数分解}
\label{b3:subsec:G01:03}

\begin{example}\label{b3:exm:real-projection-disk}
设 \(R\) 为实闭域。集合
\[
S=\{(x,y):y^2+x^2=1,\ y\ge x\}
\]
在 \(x\) 轴上的投影为 \(\{-1\le x\le0\}\cup\{0\le x,\ 2x^2\le1\}\)。事实上，圆上有点首先要求 \(-1\le x\le1\)。若 \(x\le0\)，取非负根 \(y=\sqrt{1-x^2}\) 即可；若 \(x>0\)，可取的最大 \(y\) 为该非负根，故存在满足条件的 \(y\) 当且仅当 \(1-x^2\ge x^2\)。这一计算明确使用了实闭域中非负元有平方根，而非仅使用复数域上的消去。
\end{example}

\term{柱状代数分解}[cylindrical algebraic decomposition, CAD] 把实空间分成有限个半代数胞腔，使指定多项式在每个胞腔上符号不变，并使各胞腔向低维坐标空间的投影相同或不相交。在二维，先把 \(x\) 轴按若干关键根分成点与区间；在各区间上排列 \(y\) 方向的根图像，再取图像之间的带状区域。这正是“柱状”的含义，并非说各胞腔都是直积。

\ProseParagraphBreak
上例的关键参数包括 \(\pm1\) 以及 \(2x^2-1\) 的根。一般计算还须处理首项系数消失、重根与两个多项式共有根等情形；判别式与结式提供候选投影多项式，但退化时不能只凭一份通用公式完成分解。\secref{b3:sec:G05}给出一维根计数和隔离的完整步骤；高维 CAD 的正确性与复杂性证明留待实代数几何专门课程。
